% !TEX root = ../main.tex

\chapter{收缩流道对熄火边界的定量影响}

\section{引言}

基于YJH-1E紫外火焰探测器采集的实验数据，本章利用已建立的同温工况和占空比判据，对比无收缩段和有收缩段构型的火焰稳定性差异，量化收缩段对熄火边界的提升效果。本章首先介绍实验条件与方法，然后呈现有/无收缩段的占空比对比结果，最后从主流加速效应、压力梯度效应和剪切层增强三个角度分析收缩段影响火焰稳定性的物理机制，并与Wang的Damkohler数模型进行定性对比。

\section{实验条件与方法}

实验在同温工况下开展，主流温度约 \(1342\mathrm{K}\) 。对比两种构型：无收缩段和 \(60\mathrm{mm}\) 五次曲线收缩段。每个工况稳定燃烧五分钟，记录火焰信号。在主流条件、二级射流参数相同的条件下进行对比。

\section{实验结果}

在空气流量820SLM同温工况下，无收缩段构型占空比为0.14，而加入 \(60\mathrm{mm}\) 五次曲线收缩段后，占空比提升至0.24。两者对比可以明显看出，收缩段的加入对火焰稳定性有积极影响，占空比提升约 \(71\%\)。

\begin{table}[!htbp]
	\centering
	\caption{有/无收缩段占空比对比}
	\begin{tabular}{c|c|c}
		\hline
		构型 & 占空比 & 相对于无收缩段的提升 \\
		\hline
		无收缩段 & 0.14 & — \\
		60mm五次曲线收缩段 & 0.24 & +71\% \\
		\hline
	\end{tabular}
\end{table}

后续，又对 \(60\mathrm{mm}\) 与 \(32\mathrm{mm}\) 五次曲线收缩段构型进行不同流量同温工况实验。每个工况稳定燃烧五分钟，记录火焰信号。

\begin{figure}[H]
	\centering
	\includegraphics[width=0.7\textwidth]{fig14}
	\caption{32mm占空比曲线（820 SLM工况）}
	\label{fig:32mm}
\end{figure}

\begin{figure}[H]
	\centering
	\includegraphics[width=0.7\textwidth]{fig15}
	\caption{60mm占空比曲线（820 SLM工况）}
	\label{fig:60mm}
\end{figure}

图\ref{fig:32mm}和图\ref{fig:60mm}分别展示了32mm和60mm收缩段两种构型在五个特征工况点（360、510、650、820、970 SLM）的占空比变化曲线。需要说明的是，当前选取的是流量跨度较大的特征工况点，覆盖从稳定燃烧到完全吹熄的完整区间，后续需补充中间流量工况点的数据以完善占空比曲线。

\section{分析与讨论}

综合以上数据，初步可以得出：在相同的主流工况下，加入收缩段后火焰占空比有一定提升，说明收缩段的加入对火焰稳定性有积极影响。可能源于收缩段对主流流场的改变。从物理机制分析，收缩段的存在使主流沿流向加速，增强了射流与主流之间的剪切层混合，有利于燃料与空气的掺混，从而在一定程度上延迟了吹熄的发生。而对比 \(32\mathrm{mm}\) 收缩段和 \(60\mathrm{mm}\) 收缩段数据可知， \(32\mathrm{mm}\) 收缩段在相同工况下较 \(60\mathrm{mm}\) 收缩段火焰稳定性更高，可能的原因是：收缩段使主流加速，增强了射流与主流之间的剪切层混合，有利于燃料与空气的掺混和火焰稳定；但过强的加速（ \(60\mathrm{mm}\) 收缩段）可能增大火焰面的拉伸应变率，使局部Karlovitz数超过临界值，反而削弱稳定效果。

\subsection{与Wang提出的Damkohler数模型的定性对比}

Wang等人\cite{wang2015}基于Damkohler数建立了横向射流吹熄的理论模型，将吹熄判据表述为混合时间与化学反应时间的比值低于临界值：

\begin{equation}
	Da = \frac{\tau_{m}}{\tau_{c}} = \frac{L / u}{\alpha / S_{L}}
\end{equation}

其中 \(\tau_{m}\) 为特征混合时间， \(\tau_{c}\) 为特征化学反应时间， \(L\) 为特征混合长度， \(u\) 为特征速度， \(\alpha\) 为热扩散率， \(S_{L}\) 为层流火焰传播速度。该模型表明，吹熄临界条件取决于混合过程与化学反应之间的竞争。

在本实验中，收缩段的加入改变了主流的速度场和压力场，进而影响混合时间尺度。收缩段使主流加速，一方面可能减小混合时间 \(\tau_{m}\) ，使Damkohler数降低，趋于吹熄；另一方面，增强的剪切层混合可能改变特征混合长度 \(L\) ，产生相反的效应。两种效应的竞争决定了最终的稳定效果。

本实验观测到的“ \(32\mathrm{mm}\) 收缩段优于 \(60\mathrm{mm}\) 收缩段”的趋势，可以用上述框架来解释：适度的主流加速（ \(32\mathrm{mm}\) 收缩段）增强了剪切层混合，有利于燃料与空气的掺混，使火焰能够在更宽的工况范围内维持稳定；而过强的加速（ \(60\mathrm{mm}\) 收缩段）可能使混合时间过短，导致Damkohler数低于临界值，反而促进吹熄。这一趋势与Wang模型中Damkohler数随混合时间单调变化的基本物理图像定性一致。

\section{本章小结}

本章对比了有无收缩段、不同高度收缩段的数据，结论如下：

（1）通过820SLM工况下有/无收缩段的对比实验，无收缩段构型占空比为0.14，60mm五次曲线收缩段构型占空比为0.24，占空比提升约71\%。初步表明收缩段的加入对火焰稳定性有积极影响，可能源于收缩段使主流沿流向加速，增强了射流与主流之间的剪切层混合，有利于燃料与空气的掺混。

（2）通过60mm与32mm两种收缩段构型在五个特征工况点（360、510、650、820、970SLM）的占空比对比，32mm收缩段在各近极限工况下的占空比均高于60mm收缩段。在820SLM工况下，32mm收缩段占空比为0.628，而60mm收缩段仅为0.233；在970SLM工况下，32mm收缩段仍有0.277，而60mm收缩段已降至0.000，完全吹熄。在实验涉及的两种收缩高度中，32mm收缩段效果优于60mm收缩段，初步表明存在一个最优收缩比范围。

（3）收缩段高度对火焰稳定性的影响可能源于两种效应的竞争：适度的主流加速有利于增强剪切层混合，促进火焰稳定；而过强的加速则可能从两方面产生负面影响——增大火焰面的拉伸应变率使局部Karlovitz数超限，或缩短混合时间使Damkohler数低于临界值。该趋势与Wang的Damkohler数模型定性相符。